%% ICAIF-style paper backbone for the CVaR-Constrained RL Portfolio Allocator.
%% Format intentionally mirrors recent ICAIF submissions (acmart, sigconf).
%% This is a BACKBONE: method/setup are committed; result tables are placeholders
%% (\TODO) to be filled by the pipeline. Every committed number below is drawn
%% from configs/experiment.yaml.
\documentclass[sigconf,nonacm]{acmart}

\usepackage{booktabs}
\usepackage{amsmath}
\usepackage{graphicx}
\usepackage{xcolor}
\graphicspath{{figures/}{../results/}}

%% --- Macro for "pending results" placeholders; keeps the draft honest. ---
\newcommand{\TODO}[1]{\textcolor{red}{[\textsc{todo}: #1]}}

\begin{document}

\title{CVaR-Constrained Reinforcement Learning for Benchmark-Aware
Multi-Asset Allocation}

%% Replace with the final author block before submission.
\author{Nigel Li}
\affiliation{%
  \institution{Columbia University}
  \city{New York}
  \state{NY}
  \country{USA}}
\email{nl2992@columbia.edu}

\renewcommand{\shortauthors}{Li}

\begin{abstract}
Mean--variance optimisation underpins modern portfolio construction but
penalises upside and downside symmetrically and is silent on path-wise
drawdown---precisely the risks that matter when allocating across asset classes
through stress regimes. We study a reinforcement-learning (RL) allocator that
optimises return while \emph{explicitly} controlling tail risk and drawdown,
formulated around a cost-adjusted Conditional Value-at-Risk (CVaR) objective with
hard constraints on per-period CVaR, drawdown and concentration. The allocator
acts in a compact, testable multi-asset environment (equity index, rates, credit,
commodity, FX) with weekly rebalancing, proportional transaction costs and a
long-only budget. We benchmark learned policies against equal-weight and
inverse-volatility allocation, reporting annualised return, Sharpe, Sortino,
maximum drawdown, $\mathrm{CVaR}_{95}$, turnover and constraint-violation rates,
with paired-bootstrap significance on per-period returns.
\emph{[Backbone: headline results pending the trained-agent runs.]}
On the synthetic study the CVaR-constrained allocator is expected to reduce
$\mathrm{CVaR}_{95}$ and maximum drawdown relative to the risk-agnostic baselines
at comparable or better risk-adjusted return; we will report \TODO{headline
$\mathrm{CVaR}_{95}$ / max-DD / Sharpe deltas} and ablate the contribution of the
CVaR term and the drawdown constraint.
\end{abstract}

\begin{CCSXML}
<ccs2012>
<concept><concept_id>10010147.10010257</concept_id>
<concept_desc>Computing methodologies~Reinforcement learning</concept_desc>
<concept_significance>500</concept_significance></concept>
<concept><concept_id>10003752</concept_id>
<concept_desc>Applied computing~Economics</concept_desc>
<concept_significance>300</concept_significance></concept>
</ccs2012>
\end{CCSXML}
\ccsdesc[500]{Computing methodologies~Reinforcement learning}
\ccsdesc[300]{Applied computing~Economics}

\keywords{portfolio allocation, reinforcement learning, Conditional
Value-at-Risk, drawdown control, risk-constrained optimisation, mean-variance}

\maketitle

%% =====================================================================
\section{Introduction}
%% Context -> problem -> research question -> contributions -> roadmap.

Markowitz mean--variance optimisation (MVO) laid the foundation of modern
portfolio theory~\cite{markowitz1952portfolio}, but its variance penalty treats
upside and downside symmetrically and ignores the \emph{path} of wealth: two
portfolios with identical variance can have very different drawdowns. For
multi-asset allocators that must survive stress regimes, the economically
relevant risks are \emph{tail loss} and \emph{drawdown}, not variance. A growing
literature applies machine learning to portfolio construction---both to estimate
MVO inputs better~\cite{lee2024dfl} and to recommend portfolios that respect
investor preferences~\cite{chung2023mvecf}---and reinforcement learning offers a
natural way to optimise a risk measure \emph{directly}, end to end, under
transaction costs and constraints.

We therefore ask:
\begin{quote}
\emph{Can a reinforcement-learning allocator that optimises a cost-adjusted CVaR
objective, subject to explicit drawdown and concentration constraints, control
tail risk and drawdown better than risk-agnostic baselines without sacrificing
risk-adjusted return?}
\end{quote}

\noindent\textbf{Our main contributions are:}
\begin{itemize}
\item A compact, testable \textbf{multi-asset allocation environment} with budget
normalisation, proportional transaction costs and a long-only weight cap, plus a
dataset contract for returns, factor exposures, costs and benchmark weights
(Section~\ref{sec:method}).
\item A \textbf{cost-adjusted CVaR objective with hard risk constraints}
(per-period $\mathrm{CVaR}_{95}$, maximum drawdown, per-asset concentration),
expressed so that constraint violations are measurable and reportable
(Section~\ref{sec:method}).
\item A \textbf{benchmark suite and evaluation protocol}: equal-weight and
inverse-volatility baselines, a held-out evaluation with a controlled stress
regime, and significance testing on per-period returns
(Section~\ref{sec:setup}).
\item \emph{[Pending]} An \textbf{empirical study} quantifying the tail-risk and
drawdown reduction from the CVaR objective and the drawdown constraint, with
ablations isolating each ingredient (Section~\ref{sec:results}).
\end{itemize}

The remainder reviews background and related work
(Section~\ref{sec:related}), details the environment, objective and policies
(Section~\ref{sec:method}), the experimental setup (Section~\ref{sec:setup}) and
results (Section~\ref{sec:results}), then discusses limitations, future work and
conclusions (Sections~\ref{sec:limitations}--\ref{sec:conclusion}).

%% =====================================================================
\section{Background and Related Work}
\label{sec:related}

\paragraph{Mean--variance optimisation and its inputs.} MVO~\cite{markowitz1952portfolio}
is sensitive to input estimates; decision-focused learning trains return
predictors to improve downstream MVO decisions rather than raw forecast
accuracy~\cite{lee2024dfl}, and mean--variance-efficient collaborative filtering
blends user preference with portfolio theory~\cite{chung2023mvecf}. These works
improve the \emph{inputs} or \emph{recommendation} layer of a variance-based
objective; we instead optimise a \emph{tail} risk measure directly.

\paragraph{Conditional Value-at-Risk.} CVaR is a coherent risk measure with a
convex Rockafellar--Uryasev representation~\cite{rockafellar2000optimization}
amenable to optimisation and to use as an RL reward. We adopt it both as the
training objective and as a hard per-period constraint.

\paragraph{Reinforcement learning for allocation.} RL frames allocation as a
sequential decision problem and can optimise non-differentiable, path-dependent
objectives such as drawdown directly. Our environment is deliberately small and
swappable so that policy and evaluation interfaces can later host SAC/PPO or a
continuous-time actor--critic.

%% =====================================================================
\section{Method}
\label{sec:method}

\paragraph{Allocation environment.} At each weekly step the agent outputs target
weights $w_t$ over $N$ assets; weights are clipped non-negative and budget-
normalised (long-only, $\sum_i w_{t,i}=1$, $w_{t,i}\le w_{\max}$). Transaction
costs are charged on turnover, $c_t = \kappa \lVert w_t - w_{t-1}\rVert_1$, and
wealth compounds as $W_{t+1} = W_t\,(1 + w_t^{\top} r_t - c_t)$. The committed
configuration is summarised in Table~\ref{tab:config}.

\begin{table}[t]
\centering
\caption{Committed environment and risk configuration
(\texttt{configs/experiment.yaml}).}
\label{tab:config}
\small
\begin{tabular}{ll}
\toprule
parameter & value \\
\midrule
universe          & equity index, rates, credit, commodity, FX \\
rebalance         & weekly \\
transaction cost  & 2.0 bps on turnover \\
max weight        & 0.60 (long-only) \\
$\mathrm{CVaR}$ level $\alpha$ & 0.95 \\
weekly CVaR limit & 0.035 \\
drawdown limit    & 0.15 \\
\bottomrule
\end{tabular}
\end{table}

\paragraph{Objective and constraints.} The agent maximises a cost-adjusted CVaR
utility
\begin{equation}
\mathcal{U} = \mathbb{E}[r^{p}] - \lambda\,\mathrm{CVaR}_{\alpha}(\text{loss}),
\end{equation}
in the convex Rockafellar--Uryasev form~\cite{rockafellar2000optimization},
subject to hard constraints on per-period $\mathrm{CVaR}_{95}$ ($\le 0.035$),
maximum drawdown ($\le 0.15$) and per-asset concentration ($w_{t,i}\le 0.60$).
Constraint violations are tracked and reported as a first-class metric.

\paragraph{Policies.} Baselines are \emph{equal-weight} ($w_i = 1/N$) and
\emph{inverse-volatility} ($w_i \propto 1/\hat\sigma_i$ over a trailing window).
The learned allocator shares the same environment, cost model and constraints;
the policy/training interface is built to accept a discrete-time SAC/PPO agent or
a continuous-time actor--critic approximation. \TODO{finalise the trained-agent
class for the headline run.}

%% =====================================================================
\section{Experimental Setup}
\label{sec:setup}

\paragraph{Data.} The synthetic study draws correlated weekly returns for the
five asset classes with one embedded stress regime (a multi-week drawdown
shock), so that a risk-agnostic policy is visibly punished in the tail. The same
dataset contract (returns, factor exposures, costs, benchmark weights, optional
liabilities) accepts real multi-asset panels for the empirical extension.

\paragraph{Protocol.} Train on in-sample windows; evaluate on disjoint held-out
windows including the stress regime. We report annualised return, Sharpe,
Sortino, maximum drawdown, $\mathrm{CVaR}_{95}$, turnover and constraint-
violation rate, with paired-bootstrap confidence intervals and the Wilcoxon
signed-rank test on per-period returns. Every run records seed and configuration
for reproducibility.

%% =====================================================================
\section{Results}
\label{sec:results}

\emph{This section is a backbone; the tables below are placeholders to be filled
by the pipeline.}

\begin{table}[t]
\centering
\caption{Held-out comparison (synthetic). Lower $\mathrm{CVaR}_{95}$/max-DD/
turnover is better; higher Sharpe/Sortino is better. \TODO{populate from
\texttt{results/}}.}
\label{tab:results}
\small
\begin{tabular}{lrrrrrr}
\toprule
policy & ann.\ ret. & Sharpe & Sortino & max-DD & $\mathrm{CVaR}_{95}$ & turn. \\
\midrule
equal-weight        & \TODO{} & \TODO{} & \TODO{} & \TODO{} & \TODO{} & \TODO{} \\
inverse-vol         & \TODO{} & \TODO{} & \TODO{} & \TODO{} & \TODO{} & \TODO{} \\
\textbf{CVaR-RL (ours)} & \TODO{} & \TODO{} & \TODO{} & \TODO{} & \TODO{} & \TODO{} \\
\bottomrule
\end{tabular}
\end{table}

\begin{table}[t]
\centering
\caption{Ablations (synthetic). \TODO{populate}.}
\label{tab:ablation}
\small
\begin{tabular}{lrr}
\toprule
variant & $\mathrm{CVaR}_{95}$ & max-DD \\
\midrule
full (CVaR + drawdown constraint) & \TODO{} & \TODO{} \\
objective = mean-only (no CVaR)   & \TODO{} & \TODO{} \\
no drawdown constraint            & \TODO{} & \TODO{} \\
no transaction costs              & \TODO{} & \TODO{} \\
\bottomrule
\end{tabular}
\end{table}

We expect the CVaR-constrained allocator to dominate the risk-agnostic baselines
on $\mathrm{CVaR}_{95}$ and maximum drawdown through the stress regime, and the
mean-only ablation to inflate the tail---mirroring the role the risk objective
plays in related risk-aware studies. \TODO{state the realised deltas and their
significance.}

%% =====================================================================
\section{Limitations}
\label{sec:limitations}

The headline study uses a synthetic multi-asset market with a single embedded
stress regime; real panels will introduce non-stationarity, estimation error and
richer correlation dynamics. Costs are proportional to turnover with no market
impact or borrow modelling, and the policy is long-only with a fixed weekly
cadence. The trained-agent class and the choice of absolute versus
benchmark-relative CVaR are still being finalised (Section~\ref{sec:method}).

%% =====================================================================
\section{Conclusion and Future Work}
\label{sec:conclusion}

We present a benchmark-aware, CVaR-constrained RL allocator and an evaluation
protocol that treats tail loss, drawdown and constraint violations as
first-class outcomes rather than by-products of a variance objective. Once the
trained-agent runs complete, we will quantify the tail-risk and drawdown
reduction over risk-agnostic baselines and ablate each ingredient. Future work
includes a discrete-time SAC/PPO and a continuous-time actor--critic agent,
benchmark-relative and surplus-shortfall risk targets, leverage and shorting
variants, and a real multi-asset empirical study spanning historical stress
episodes.

\begin{acks}
Artefacts are produced by the open pipeline in the project repository and are
reproducible from a clean checkout.
\end{acks}

\bibliographystyle{ACM-Reference-Format}
\bibliography{references}

\end{document}
